% critique.tex — f164v, vv. 165–208 : commentaire critique
% Articles textuels tirés du commentaire de l'auteur (latex/Aristophane-Thesmo Harmattan 2023 vv 1-641.tex),
% texte et notes conservés ; répartition détaillée dans repartition.md.
\cc{176}{λέγε.}{\latin{Nam neque Euripides neque Agatho Mnesilochi jocos solent curare. 
Apparet vero, orationem sic cohaerere},\citComm{\grec{παῦσαι βαΰζων — ἀλλ᾽ ... ἔα μ᾽ εἰπεῖν}\\
\latin{desine me maledictis insequi, sed potius ea me patere loqui.}}\latin{Jam necessario is, qui rogatus est, respondere debet, estque certo verbum λέγε dic igitur Agathonis.} (Note de Fritzsche.)}
\cc{179}{ἐγὼ δὲ κοινῇ ξυμφορᾷ πεπληγμένος.}{L'émendation de Biset : substituer καινῇ à κοινῇ est reprise par quasiment tous les éditeurs. Elle a l'avantage de concorder avec la scholie : μεγάλῃ καὶ θαυμαστῇ (καινός désigne au propre ce qui est \italique{nouveau}, différent de ce qui s'est produit jusqu'à présent, et par suite, ce qui n'était pas attendu, \italique{étrange}, \italique{extraordinaire} – avec une évolution toute identique à notre \italique{in-ouï}). 
Nous sommes devant l'un des cas\nc{Voir pour un autre exemple le commentaire au vers 1174.} où la scholie rend compte d'un texte en meilleur état que celui auquel elle est attenante. (À noter qu'Euripide emploie l'expression καινάς συμφοράς aux vers 1226-7 de l'\italique{Oreste}.)\\
Cette glose s'est retrouvée dans la \italique{Souda} avec la mauvaise vedette (à l'article Κοινολογία). Pis : Suidas -- ou l'un des rédacteurs de la \italique{Souda} si l'on préfère --, voyant le problème aura ajouté aux deux termes dont le sens est incompatible avec κοινῇ un tiers, qui, lui, l'est, ἀνθρωπίνῃ.
Il écrit : κοινῇ δὲ ἀντὶ τοῦ ἀνθρωπίνῃ, μεγάλῃ καὶ θαυμαστῇ... Ce sont là des méthodes bien peu rigoureuses qui doivent nous le faire tenir en suspicion\nc{Voir \italique{infra} le commentaire au vers 735 pour une erreur analogue de Suidas.}.}
\cc{204}{νυκτερείσια.}{La forme canonique est νυκτερήσια, \italique{nocturne} (attestée tardivement par Lucien et Sextus Empiricus), forgée en miroir de ἡμερήσιος. Gardons-nous pour autant de retoucher la graphie du codex de Ravenne ; la réplique perdrait de son sel aristophanesque : νυκτ-ερείσια fait entendre le substantif du verbe ἐρείδειν, « \latin{quod in comoedia est, subagitare mulierem} » (Fritzsche).}
